\section{Experiments}
\label{sec:experiments}

The fixed-support control first isolates allocation from support. We then ask
whether the same intervention remains consequential when rotary channels are
scarce, when range scaling is applied, and when models are obtained by
continuation or adaptation. Each protocol retains its own metric and control.

%% =========================================================================
\subsection{Identifying allocation at fixed support}
\label{sec:exp-identify}

\paragraph{Protocol.}
Within each of three $151.9$M training seeds, the two arms share architecture,
trainable initialisation, token order, optimiser, LR schedule, global batch, a
$499{,}974{,}144$-token budget, and $32$ frozen evaluation anchors. The
baseline implements the uniform-in-log rule of
\citet[\S6.1]{oka2026fmrope} at training base $256$; the intervention uses
anchored \evq{} with $\tau{=}4$, obtained by normalising its quantiles to the
FMRoPE endpoints. Because their sampled extrema
and log-frequency span coincide, only the $K{-}2=30$ interior frequencies
move. Support normalisation also preserves the deployed Cosh table's $z$ in
Eq.~\eqref{eq:table-decomposition}.

\paragraph{Result.}
Every training seed favours anchored \evq{} at every OOD length. The
anchored \evq{} minus FMRoPE NLL difference is $+0.026$ at the
$256$-token training length and
$-0.281/-0.176/-0.146$ at $512$/$1$K/$2$K
(Fig.~\ref{fig:evidence-overview}a). No geometric base can reproduce this intervention
once the sampled support is pinned: \emph{fixed support $+$ different $z$
$\Rightarrow$ different trained-model behaviour.}

\paragraph{Cross-configuration and shape breadth.}
A pre-specified $50.9$M factorial then spans two bases, two training lengths,
three head dimensions, and three paired seeds. Across its $12$ structural
configurations, anchored \evq{} at $1.25\times$ rule strength and a deformation-matched exponential favour
non-uniform allocation in $10/12$ and $9/12$ cases
(App.~Table~\ref{tab:m4}). The $151.9$M control supplies the replicated direct
effect; the matched exponential shows that the direction is not tied to \evq{}
alone. Full results appear in App.~\ref{sec:identification-details}.

%% =========================================================================
\subsection{A finite budget survives architecture, range scaling, and scale}
\label{sec:exp-mature}

\paragraph{Scarce-channel MLA.}
A $432$M MLA transformer trains from scratch over three seeds with only
$d_{\mathrm{rope}}{=}32$, hence $K{=}16$ rotary frequencies. Geo/\evq{} PPL
is $35.4/35.8$ at the $8$K training length. At $16$K, however, it is
$138.8/95.6$: a $31.1\%$ reduction, with every seed favouring \evq{}
(Fig.~\ref{fig:evidence-overview}b; App.~Table~\ref{tab:mla}). The result
directly links the finite-$K$ view to a scarce-budget stress test.

\paragraph{Range composition.}
Allocation remains visible after applying the same range operator. In a $454$M
three-seed MHA study, a fixed \rs{} operator raises $8$K teacher-forced
retrieval from $41\%$ to $61\%$ on Geo and from $53\%$ to $100\%$ on \evq{}
(App.~Table~\ref{tab:evq-ramp}). The operator is identical; its leverage depends
on the table that shaped the weights during training.

\paragraph{Continuation and training-scale trend.}
The crossover also appears after full-parameter continuation. Starting both
$750$M arms from a shared $2$K Geo checkpoint and continuing all parameters for
$500$M tokens at $4$K gives Geo/\evq{} PPL $22.0/22.3$ at $4$K but
$45.1/24.4$ at $16$K; strict autoregressive exact retrieval at $8$K changes
from $0\%$ to $77.5\%$ (App.~Table~\ref{tab:750m}). At $1.485$B, a
same-initialisation, same-scientific-recipe from-initialisation comparison with
a $2.097$B-token budget crosses in favour of \evq{} at $8$K and $16$K;
$122/128$ and $126/128$ held-out documents favour it at those lengths. Exact
PPL values and protocol details appear in App.~\ref{sec:mature-details}. This is
the end of the full-parameter training-scale trend; the $8$B evidence below is
adaptation.

%% =========================================================================
\subsection{From effective context to remote-content use}
\label{sec:exp-capability}

\paragraph{$1.485$B real-document QA.}
Starting from the two OLMo-2 Stage-A parents, both arms freeze inherited V/O
LoRA and continue only Q/K for $300$ steps with identical data, order,
optimiser, trainable count, and $4$K/$8$K/$16$K phase exposure; physical
sequences remain at most $4$K. On held-out 2WikiMultiHopQA, Native/\evq{}
token-F1 is $25.99/24.84\%$ at $4$K, $0.07/21.48\%$ at $8$K, and
$0/8.57\%$ at $16$K (Fig.~\ref{fig:evidence-overview}c)
\citep{bai2024longbench}. Thus Q/K adaptation keeps the in-window QA score
close to Native while only the \evq{} arm carries substantial answer overlap
to $8$K. An
independent continuation over all $13$ RULER families shows the same $2\times$
direction, with official macro $2.02/31.63\%$ at $8$K
\citep{hsieh2024ruler}; all lengths and exact-match values appear in
App.~\ref{sec:mature-details}.

\paragraph{$8$B causal remote-source use.}
On LLaMA-3-8B-Instruct \citep{grattafiori2024llama3}, rank-$64$ Q/K/V/O LoRA
uses the same adaptation pipeline for Native and \evq{}. PPL is $6.82/10.07$,
$108.96/24.07$, and $991.48/127.91$ at $8$/$16$/$32$K, with the long-range
direction holding on all $24$ frozen natural-text packs and all three source
domains. More importantly, at true $16$K, median target-block hit@16 rises
from $18.75\%$ to $64.06\%$, and deleting the remote gold block changes NLL by
$-0.0095/+1.5055$. The deletion intervention shows that the improved
long-position probability depends on the remote source rather than a local
readout artifact. The task-family continuation is reported in
App.~\ref{sec:mature-details}.

A two-seed $129.6$M video DiT also favours \evq{} on training, all-frame
extrapolation, and far-frame MSE in both runs
(App.~\ref{sec:video-dit}), extending the evidence to bidirectional 3D RoPE.

%% =========================================================================
\subsection{Beyond a single formula}
\label{sec:exp-tau}

Two controls separate the allocation effect from the particular \evq{} formula.
A deformation-matched exponential improves in $9/12$ factorial configurations,
while the three pre-specified \evq{} multipliers ($0.75\times$, $1.00\times$,
$1.25\times$) all reduce aggregate weighted OOD NLL relative to Geo. Their
ordering varies by configuration, so $c{=}1$ is used as a zero-search operating
prior rather than a tuned point optimum. Non-uniform allocation supplies the
empirical direction, and \evq{} supplies a closed-form member of that design
space. Full multiplier, boundary, and matched-shape results appear in
App.~\ref{sec:identification-details}.
